\documentclass[11pt]{article}
\input{../calcnotes.sty}
\fancyhead[C]{1.1 Lines}

\begin{document}

\noindent Calculus is a different type of math. In most algebra classes, we talk about exact amounts, or inequalities that are exact. Calculus deals with quantities that are changing, rates of change, and increments. The increment is defined as

\[\Delta x = x_2 - x_1 \text{ and } \Delta y = y_2 - y_1\]

\noindent This looks just like the slope of a line:\\[2ex]

\noindent Which often leads us straight to point-slope form:\\[2ex]

\noindent Which often leads us straight to slope-intercept form:\\[2ex]

\noindent Which sometimes leads us to the General Linear Equation:\\[2ex]

\noindent (Did we always call it that?)\\[2ex]

\noindent Find the slope of a line through the points $\displaystyle (3,-2)$ and $\displaystyle (4,1)$. Write the equation of a line in all three forms.\\[2ex]

\noindent What about parallel and perpendicular lines? Write an equation for a line through the point $\displaystyle (-1,2)$ that is parallel and perpendicular to the line $\displaystyle y=3x-4$.\\[2ex]

\noindent What line is perpendicular to the line $\displaystyle x=-3$ and goes through the point $\displaystyle (-3,4)$?\\[2ex]

\noindent How would we find the equation of a line given the data below (and what is the $\displaystyle f(x)$ notation)?\\[1ex]

\begin{tabular}{ |c|c| }
    \hline
    $\displaystyle x$ & $\displaystyle f(x)$          \\
    \hline
    -1                & $\displaystyle \frac{14}{3}$  \\
    1                 & $\displaystyle \frac{-4}{3}$  \\
    2                 & $\displaystyle \frac{-13}{3}$ \\
    \hline
\end{tabular}\\[2ex]

\noindent Try it again with Fahrenheit and Celsius (what temperatures do you already know about these two scales?).\\[2ex]

\newpage

\noindent And finally, how do you model a whole set of data? By hand? By calculator? Find a linear model for the world population shown.\\[1ex]
\begin{tabular}{|c|c|}
    \hline
    \multicolumn{2}{|c|}{\textbf{World Population}} \\ \hline
    \textbf{Year} & \textbf{Population (millions)}  \\ \hline
    1980          & 4454                            \\ \hline
    1985          & 4853                            \\ \hline
    1990          & 5285                            \\ \hline
    1995          & 5696                            \\ \hline
    2003          & 6305                            \\ \hline
    2004          & 6378                            \\ \hline
    2005          & 6450                            \\ \hline
\end{tabular}

\end{document}
